% !TeX root = ../../Book1.tex
% 纲要标题：### 11.3 主理想整环
% 纲要位置：计划纲要.md，第 428 行。

\section{主理想整环}
\label{b1:sec:1103}

% 纲要内容（仅作写作计划，尚非教材正文）：
%
% * 理想主性与素元性质
% * 主理想整环是唯一分解整环
% * 非 Euclid 的主理想整环作为选读例子
%

Euclid 算法提供了求最大公因子的方法，但它还有一个不依赖具体计算的后果：每个理想都能用一个元素生成。把这一结论单独提出，许多分解论证明便不再需要一个指定的除法算法。最后的例子将说明，这样做确实扩大了研究范围。

\subsection{理想主性与素元性质}
\label{b1:subsec:1103:01}

\begin{definition}\label{b1:def:pid}
设 \(R\) 是交换整环。若对每个理想 \(I\subseteq R\)，都存在 \(d\in R\) 使 \(I=(d)=\{\,dr\mid r\in R\,\}\)，则称 \(R\) 为\term[zhulixiangzhenghuan]{主理想整环}[principal ideal domain]，简称 PID。
\end{definition}

定义包含零理想 \((0)\) 和单位理想 \((1)\)。域只有这两个理想，因而是主理想整环；“整环”这一要求不能删去后仍沿用同一个概念。

\ProseParagraphBreak
要从 Euclid 除法得到理想主性，可以把除法用在理想内部：选定一个非零元素 \(d\) 后，任何不被 \(d\) 整除的理想元素都会产生一个非零余数，而这个余数仍在理想中。因此，若 \(d\) 的 Euclid 值已经最小，就不能再有这样的元素，\(d\) 必须整除整个理想。

\begin{theorem}\label{b1:thm:euclidean-pid}
设 \(R\) 为 Euclid 整环，则 \(R\) 是主理想整环。
\end{theorem}
\begin{proof}
设 \(R\) 带 Euclid 函数 \(\delta\)，\(I\) 是非零理想。从 \(I\setminus\{0\}\) 中取 \(d\)，使 \(\delta(d)\) 最小。对任意 \(a\in I\)，写 \(a=qd+r\)，其中 \(r=0\) 或 \(\delta(r)<\delta(d)\)。但 \(r=a-qd\in I\)，最小性排除了非零余数，所以 \(a\in(d)\)。反向包含由 \(d\in I\) 得出，故 \(I=(d)\)。零理想本来就由 \(0\) 生成。
\end{proof}

\begin{proposition}\label{b1:prop:pid-bezout}
设 \(R\) 为主理想整环，\(a,b,d\in R\)。则 \(a,b\) 有最大公因子，且
\[
(a,b)=(d)\quad\Longleftrightarrow\quad d\text{ 是 }a,b\text{ 的最大公因子}.
\]
此时存在 \(u,v\in R\) 使 \(d=ua+vb\)。
\end{proposition}
\begin{proof}
主性给出 \((a,b)=(d_0)\)。其中 \(a,b\in(d_0)\) 表明 \(d_0\) 是公因子，而 \(d_0\in(a,b)\) 给出线性组合表达式；\cref{b1:prop:bezout-gcd}于是说明它是最大公因子。任一其他最大公因子 \(d\) 都与 \(d_0\) 相伴，故\cref{b1:prop:associate-ideals}给出 \((d)=(d_0)=(a,b)\)，再由 \(d\in(a,b)\) 得到相应 Bézout 系数。这也包含 \(a=b=0\) 的情形。
\end{proof}

\begin{proposition}\label{b1:prop:pid-irreducible-prime}
在主理想整环 \(R\) 中，每个不可约元 \(p\) 都是素元，且 \((p)\) 是极大理想。因此每个非零素理想都是极大理想。
\end{proposition}
\begin{proof}
设 \((p)\subseteq I\subseteq R\)。写 \(I=(d)\)，则 \(p=dc\)。不可约性使 \(d\) 为单位或 \(c\) 为单位；前者给 \(I=R\)，后者给 \(I=(p)\)。又 \(p\) 非单位，所以 \((p)\) 是极大理想。由\cref{b1:prop:prime-max-quotient}，\(R/(p)\) 是域，因而是整环，\((p)\) 为素理想；再用\cref{b1:prop:prime-implies-irreducible}的素理想判别得 \(p\) 为素元。

若 \(P\ne(0)\) 是素理想，写 \(P=(p)\)。因 \(P\) 非零且真，\(p\) 非零非单位；同一判别说明 \(p\) 为素元，再由\cref{b1:prop:prime-implies-irreducible}得它不可约。刚才的论证便说明 \(P\) 极大。
\end{proof}

\subsecref{b1:subsec:1101:03}已把不可约性解释为“在真主理想中极大”。主理想整环中每个理想都为主理想，所以这一极大性涵盖了全部真理想。上述证明中的推理可写为
\[
p\text{ 不可约}
\quad\Longrightarrow\quad
(p)\text{ 为极大理想}
\quad\Longrightarrow\quad
(p)\text{ 为素理想}
\quad\Longrightarrow\quad
p\text{ 为素元}.
\]
“每个非零素理想都是极大理想”中的非零条件也有作用。在整环中，\((0)\) 总是素理想；但在主理想整环 \(\mathbb Z\) 中，有 \((0)\subsetneq(2)\subsetneq\mathbb Z\)，所以 \((0)\) 不是极大理想。

\ProseParagraphBreak
例如对域 \(F\)，\cref{b1:prop:field-poly-division}与\cref{b1:thm:euclidean-pid}给出 \(F[x]\) 是主理想整环。若 \(p(x)\) 不可约，商 \(F[x]/(p)\) 因而是域。这个结论会在构造方程根所在的域时发挥作用；它所依赖的是理想极大性，不是预先知道根的存在。

\subsection{主理想整环的唯一分解性}
\label{b1:subsec:1103:02}

理想主性还将保证元素的不可约分解存在且唯一。这里先定义唯一分解整环，以陈述主理想整环的这一结构性质；\secref{b1:sec:1104}将脱离主理想性，单独研究唯一分解本身及其判别。

\begin{definition}\label{b1:def:ufd}
设 \(R\) 是交换整环。若每个非零非单位都能写成有限个不可约元的乘积，并且同一非零非单位 \(a\in R\) 的任意两种表达式
\[
a=u p_1\cdots p_m=v q_1\cdots q_n
\]
中，\(u,v\) 为单位、\(p_i,q_j\) 为不可约元时，必有 \(m=n\)，且可重新排列 \(q_j\)，使每个 \(p_i\) 与 \(q_i\) 相伴，则称 \(R\) 为\term[weiyifenjiezhenghuan]{唯一分解整环}[unique factorization domain]，简称 UFD。
\end{definition}

单位允许看成空乘积再乘一个单位；零不属于分解定理的对象。唯一性既不固定因子的次序，也不固定各因子相伴类中的代表。因此 \(6=2\cdot3=(-2)(-3)\) 并不是两种不同的素因数分解。

\ProseParagraphBreak
一般理想的并未必是理想，甚至未必对加法封闭。例如 \(2\mathbb Z\cup3\mathbb Z\) 包含 \(2,3\)，却不包含 \(2+3=5\)。理想升链的并有所不同：属于不同项的两个元素总能一起放入其中较大的一项，这使理想的封闭条件得以保留。

\begin{lemma}\label{b1:lem:pid-ascending-chain}
设 \(R\) 为主理想整环。则 \(R\) 的每条理想升链
\(I_1\subseteq I_2\subseteq\cdots\) 都从某项开始恒定。
\end{lemma}
\begin{proof}
令 \(I=\bigcup_{j\geq1}I_j\)。若 \(x\in I_r,y\in I_s\)，升链条件使两者都属于 \(I_{\max(r,s)}\)，故 \(x-y\) 仍在 \(I\)；任一元素乘环中元素也留在原来所属的 \(I_j\)。因此 \(I\) 是理想，写成 \((d)\)。生成元 \(d\) 属于某个 \(I_m\)，所以 \(I=(d)\subseteq I_m\subseteq I\)，从而 \(I_m=I\)，所有 \(j\geq m\) 也满足 \(I_j=I\)。
\end{proof}

分解元素与扩大主理想正好方向相反。若非零元素满足 \(a_1=a_2d_1\)，且 \(d_1\) 非单位，则 \((a_1)\subsetneq(a_2)\)。因此，若一个元素总能继续作非平凡分解，就会产生严格上升的主理想链。下面用升链终止来排除无限继续分解的情形，再用不可约元的素性比较两种分解。

\begin{theorem}\label{b1:thm:pid-ufd}
设 \(R\) 为主理想整环，则 \(R\) 是唯一分解整环。
\end{theorem}
\begin{proof}
先证分解存在。若存在不能分解成有限个不可约元乘积的非零非单位 \(a_1\)，则 \(a_1\) 本身不可为不可约元，故 \(a_1=b_1c_1\)，其中两个因子都非单位。至少有一个因子仍不能分解，否则合并两边分解便得到 \(a_1\) 的分解。把这个因子记为 \(a_2\)，另一个记为 \(d_1\)，于是 \(a_1=a_2d_1\)。有 \((a_1)\subsetneq(a_2)\)：若相等，\cref{b1:prop:associate-ideals}使 \(a_1,a_2\) 相伴，再由非零消去律迫使 \(d_1\) 为单位，矛盾。重复选择会得到严格升链
\((a_1)\subsetneq(a_2)\subsetneq\cdots\)，与\cref{b1:lem:pid-ascending-chain}矛盾。因此分解存在。

再证唯一性。设 \(u p_1\cdots p_m=v q_1\cdots q_n\)。由\cref{b1:prop:pid-irreducible-prime}，不可约元 \(p_1\) 是素元；它不整除单位 \(v\)，对乘积逐次应用素性，必整除某个 \(q_j\)。不可约性使 \(q_j=p_1w\)，其中 \(w\) 为单位。把这个因子移至首位，消去非零的 \(p_1\)，得到两个较短分解的等式，单位因子相应并入 \(v\)。如此归纳，配对因子都相伴。若一边先用尽而另一边仍有不可约因子，就会使若干非单位之积成为单位；但乘积有逆元时，每个因子也有逆元，矛盾。因此两边同时用尽，\(m=n\)，唯一性得证。
\end{proof}

这一证明的两部分使用不同的条件：
\[
\begin{array}{c@{\qquad}c}
\text{分解存在性}&\text{分解唯一性}\\
\text{理想升链终止}&\text{不可约元都是素元}
\end{array}
\]
由此前的具体 Euclid 函数，\(\mathbb Z\)、\(F[x]\) 和 \(\mathbb Z[i]\) 均为唯一分解整环。不过，主性证明中的最小值选择只断言某个生成元存在；若没有可执行的除法步骤，它本身还不是计算该生成元的算法。

\OptionalSubsection{非 Euclid 的主理想整环}{b1:subsec:1103:03}

要说明“Euclid 整环 \(\Rightarrow\) 主理想整环”的逆命题不成立，需要对同一个环完成两件事：先证明每个理想都是主理想，再排除所有可能的 Euclid 函数。前一部分将利用格点估计构造允许倍乘的较弱余数；后一部分则从任意 Euclid 函数推出一个不可能的有限商环。设
\[
\omega=\frac{1+\sqrt{-19}}2,\qquad
R=\mathbb Z[\omega]\subseteq\mathbb C.
\]
因为 \(\omega^2=\omega-5\)，每个元素唯一写成 \(a+b\omega\)，\(a,b\in\mathbb Z\)。定义
\[
N(a+b\omega)=\lvert a+b\omega\rvert^2
=a^2+ab+5b^2
=\Bigl(a+\frac b2\Bigr)^2+\frac{19b^2}4.
\]
范数是非负整数且乘法相容，非零元素的范数是正整数。若 \(u\) 是单位，则 \(N(u)N(u^{-1})=N(1)=1\)，所以 \(N(u)=1\)。由最后一个表达式，范数等于 \(1\) 时只能有 \(b=0,a=\pm1\)，而 \(\pm1\) 确实可逆，因此 \(R\) 的单位恰为 \(\pm1\)。

\ProseParagraphBreak
比较下面两种余数要求，其中 \(\beta\ne0\)。左边是以范数为度量的普通带余除法，右边只针对 \(\alpha\notin(\beta)\)：
\[
\begin{array}{c@{\qquad}c}
\text{普通带余除法}&\text{允许倍乘的较弱条件}\\
\alpha=d\beta+r&c\alpha=d\beta+r\\
r=0\ \text{或}\ N(r)<N(\beta)&0<N(r)<N(\beta)
\end{array}
\]
右边的乘数 \(c\in R\) 可以依赖于 \(\alpha/\beta\)，所以所得 \(r\) 是 \(c\alpha\) 的余数，而不是 \(\alpha\) 的余数。这一条件不提供普通带余除法，却足以证明理想主性：若非零理想 \(I\) 中的 \(\beta\) 具有最小正范数，而 \(\alpha\in I\setminus(\beta)\)，则 \(r=c\alpha-d\beta\) 仍在 \(I\) 中，其正范数严格更小，产生矛盾。令 \(z=\alpha/\beta\notin R\)，右边的不等式就化为 \(0<\lvert cz-d\rvert^2<1\)。下面的格点引理构造这样的 \(c,d\)，并保证余数非零。

\begin{lemma}\label{b1:lem:nineteen-small-combination}
令
\[
\omega=\frac{1+\sqrt{-19}}2,\qquad R=\mathbb Z[\omega]\subseteq\mathbb C.
\]
对任意 \(z\in\mathbb C\setminus R\)，存在 \(c,d\in R\)，使
\(0<\lvert cz-d\rvert^2<1\)。
\end{lemma}
\begin{proof}
由 \(R=\mathbb Z+\mathbb Z\omega\) 及 \(\operatorname{Im}\omega=\sqrt{19}/2\)，减去适当整数倍的 \(\omega\) 可以把虚部移入一个长度为 \(\sqrt{19}/2\)、以零为中心的区间；再减去整数，把实部移入长度为 \(1\)、以零为中心的区间。因此可取 \(\rho\in R\)，使调整后的 \(z_0=z-\rho=x+yi\) 落入基本矩形
\[
\lvert x\rvert\leq\frac12,\qquad
\lvert y\rvert\leq\frac{\sqrt{19}}4.
\]
由于 \(z\notin R\)，调整后的 \(z_0\) 仍不在 \(R\) 中，否则 \(z=z_0+\rho\in R\)。特别地 \(z_0\ne0\)。若 \(\lvert z_0\rvert<1\)，取 \(c=1,d=\rho\)，就有 \(0<\lvert cz-d\rvert^2=\lvert z_0\rvert^2<1\)，已经得到结论。

\cref{b1:fig:nineteen-lattice}中矩形与单位圆之间的上下两片区域，正是基本矩形内仍须处理 \(\lvert z_0\rvert\geq1\) 的点。图中取 \(z_0=\omega/2\)：从 \(z_0\) 到 \(\omega\) 是倍乘，从 \(\omega\) 回到原点才是减去格点 \(\omega\)。

\begin{figure}[htbp]
\centering
\begin{tikzpicture}[x=1.55cm,y=1.55cm,>=stealth]
  \def\h{1.08972474} % sqrt(19)/4
  \def\s{0.86602540} % sqrt(3)/2
  \fill[color=AlgebraBlueLight] (-.5,\s)--(-.5,\h)--(.5,\h)--(.5,\s)
    arc[start angle=60,end angle=120,radius=1]--cycle;
  \fill[color=AlgebraBlueLight] (-.5,-\s)--(-.5,-\h)--(.5,-\h)--(.5,-\s)
    arc[start angle=-60,end angle=-120,radius=1]--cycle;
  \draw[->,color=AlgebraGray] (-1.2,0)--(1.25,0) node[right] {\(\operatorname{Re}z\)};
  \draw[->,color=AlgebraGray] (0,-1.26)--(0,2.51) node[above] {\(\operatorname{Im}z\)};
  \draw[dashed,color=AlgebraGray] (0,0) circle (1);
  \draw[thick,color=AlgebraBlue] (-.5,-\h) rectangle (.5,\h);
  \fill (0,0) circle (1.2pt) node[below left] {\(0\)};
  \fill[color=AlgebraBlue] (.25,\h) circle (1.6pt);
  \node[anchor=east,color=AlgebraBlue] at (.17,1.30) {\(z_0=\omega/2\)};
  \fill[color=AlgebraBlue] (.5,2*\h) circle (1.6pt);
  \node[anchor=west,color=AlgebraBlue] at (.55,2*\h) {\(2z_0=\omega\)};
  \draw[->,thick,color=AlgebraBlue] (.28,\h+.08)--(.48,2*\h-.09)
    node[pos=.53,left] {\(\times2\)};
  \draw[->,thick,color=AlgebraGray] (.55,2*\h-.03)
    .. controls (1.68,2.02) and (1.72,.38) .. (.06,.04);
  \node[color=AlgebraGray] at (1.38,1.16) {\(-\omega\)};
\end{tikzpicture}
\caption[基本矩形与单位圆]{格点引理中的基本矩形与单位圆（虚线）。阴影为矩形内满足 \(\lvert z_0\rvert\geq1\) 的上下两片区域；箭头分别表示倍乘与减去格点，图中展示的是零余数实例。}
\label{b1:fig:nineteen-lattice}
\end{figure}

以下处理基本矩形内仍有 \(\lvert z_0\rvert\geq1\) 的情形。由 \(x^2\leq1/4\)，有 \(\lvert y\rvert\geq\sqrt3/2\)。对 \(2z_0\) 再作格点调整：减去虚部与 \(y\) 同号的 \(\omega\) 或 \(-\omega\)，再减去整数使实部绝对值不超过 \(1/2\)，得到 \(w=2z_0-d_0\)，且
\[
\lvert\operatorname{Im}w\rvert
\leq\frac{\sqrt{19}}2-\sqrt3<\frac12,
\qquad \lvert w\rvert^2<\frac12<1.
\]
其中严格不等式可由 \(\sqrt{19}<22/5\)、\(\sqrt3>17/10\) 核对。若 \(w\ne0\)，取 \(c=2,d=d_0+2\rho\)，便有 \(cz-d=2z_0-d_0=w\)，即得结论。

最后处理 \(w=0\)。此时倍乘 \(c=2\) 恰好把 \(z_0\) 送到格点，所得余数是零，不能用它反驳最小正范数。我们改换乘数 \(c\)，目标是制造余数 \(1/2\)，从而同时保证余数非零与模平方小于 \(1\)。令 \(\alpha=2z_0\in R\)；由于 \(z_0\notin R\)，\(\alpha\notin2R\)。

按 \(a+b\omega\) 的两个整数系数的奇偶性，商环 \(R/2R\) 恰有四个剩余类，分别由 \(0,1,\omega,1+\omega\) 代表。三个非零类的逆元分别由 \(1,1+\omega,\omega\) 代表，因为
\(\omega(1+\omega)=2\omega-5\equiv1\pmod{2R}\)。这说明 \(R/2R\) 是四元素域；当前所需的正是非零类可逆。\(\alpha\) 的类非零，故可取 \(c\in R\) 及 \(d_1\in R\)，使 \(c\alpha=1+2d_1\)。于是 \(cz_0-d_1=1/2\)，其模平方为 \(1/4\)。取 \(d=d_1+c\rho\)，便有 \(cz-d=cz_0-d_1=1/2\)，就得到原命题。

例如 \(z_0=\omega/2\) 位于图中的上方阴影区域，因为 \(\lvert z_0\rvert^2=5/4>1\)，但 \(2z_0-\omega=0\)。这时可直接取 \(c=1+\omega\)、\(d_1=\omega-3\)；由 \(\omega^2=\omega-5\) 得
\[
(1+\omega)\frac{\omega}{2}-(\omega-3)=\frac12.
\]
所以零余数分支也确实产生非零且模平方小于 \(1\) 的线性组合。
\end{proof}

\begin{proposition}\label{b1:prop:noneuclidean-pid}
令 \(\omega=(1+\sqrt{-19})/2\)，并令 \(R=\mathbb Z[\omega]\subseteq\mathbb C\)。则 \(R\) 是主理想整环，但不存在任何使 \(R\) 成为 Euclid 整环的 Euclid 函数。
\end{proposition}
\begin{proof}
设 \(I\ne(0)\) 为理想，取 \(\beta\in I\setminus\{0\}\)，使正整数 \(N(\beta)\) 最小。若某个 \(\alpha\in I\) 不属于 \((\beta)\)，则 \(z=\alpha/\beta\notin R\)。由\cref{b1:lem:nineteen-small-combination}取 \(c,d\in R\)，使 \(0<\lvert cz-d\rvert^2<1\)。于是非零元素 \(c\alpha-d\beta\in I\) 满足
\[
0<N(c\alpha-d\beta)
=N(\beta)\lvert cz-d\rvert^2<N(\beta),
\]
与最小性矛盾。所以 \(I=(\beta)\)，零理想也为主理想，得证主性。

现在排除任意 Euclid 函数。由于 \(N(2)=4\ne1\)，\(2\) 是非零非单位，因此所有非零非单位组成的集合非空。假设存在 Euclid 函数 \(\delta\)，可在这个集合中取 \(b\)，使 \(\delta(b)\) 最小。对任意 \(a\in R\) 作除法 \(a=qb+r\)；若 \(r\ne0\)，则 \(\delta(r)<\delta(b)\)。最小性使 \(r\) 不可能是非单位，因而只能是单位，即 \(1\) 或 \(-1\)。

因此每个模 \((b)\) 的剩余类都由 \(0,1,-1\) 之一代表，\(R/(b)\) 至多有三个元素。另一方面，\(b\) 非单位意味着 \(1\notin(b)\)，所以商环中的 \(1\ne0\)，至少有两个元素。合起来得到
\[
2\leq\lvert R/(b)\rvert\leq3.
\]
故它的大小只能为 \(2\) 或 \(3\)。

令 \(S=R/(b)\)，并令 \(p=\lvert S\rvert\in\{2,3\}\)。由\cref{b1:thm:lagrange}，\(1_S\ne0_S\) 在加法群中的阶是 \(p\)，所以 \(1_S\) 生成整个加法群。映射
\[
\mathbb Z/p\mathbb Z\longrightarrow S,\qquad
[n]\longmapsto n1_S
\]
于是良定义且为双射；它保持加法、乘法和单位，其中乘法相容性来自 \((m1_S)(n1_S)=(mn)1_S\)。因此它是含幺环同构，即 \(S\cong\mathbb F_p\)。

然而 \(\omega\) 的像必须满足 \(t^2-t+5=0\)。模 \(2\) 时，\(t=0,1\) 的值都是 \(1\)；模 \(3\) 时，\(t=0,1,2\) 的值依次为 \(2,2,1\)。两种商中都不可能有这样的元素，矛盾。故不存在 \(\delta\)。
\end{proof}

范数本身不能作 Euclid 函数，还不足以推出环不是 Euclid 整环；必须排除所有可能的 Euclid 函数。上述证明的后半部分正是通过有限商环的必要条件完成这一点。主性证明使用的乘数 \(c\) 则足以产生理想中的更小元素，却不提供对原元素的普通带余除法。

\ProseParagraphBreak
关于这个例子，可参见 Keith Conrad 的讲义\cite[Theorems 5.18, 5.22]{ConradRemarksEuclideanDomains}：定理 5.18 提炼了非域 Euclid 整环中某个非单位模下每个剩余类可由零或单位代表的必要条件；定理 5.22 讨论本小节的二次环，其中主理想性只给出证明提纲，可与此处的格点引理作方法对照。
